\chapter{Introduction}

\section{Motivations}

The idea for this project came directly from discussions within our team. Arihant wanted to work on distributed systems, but the specific project proposal and application concept were developed by Su and Nizar based on real problems they observed. 

Specifically, Nizar observed this difficulty during EPITA's \emph{Challenge Week}. In a high-pressure team setting, group members who were unfamiliar with Git spent a lot of time troubleshooting setup issues, broken repositories, and merge conflicts. These technical hurdles regularly delayed collaboration, proving that for non-programmers or short-term work, Git is too complex to be practical.

At the same time, we saw the limits of simpler sharing methods. During a database course last semester, three of our team members ended up with four different versions of an ER diagram file. Reconciling those changes manually across email, WhatsApp, and USB drives was error-prone, and the final submission used an outdated version. 

Existing tools address parts of this problem. Google Drive \cite{googledrive} and Dropbox \cite{dropbox} synchronize files through cloud servers but require internet access, accounts, and have storage quotas. Git \cite{gitbook} tracks changes perfectly but requires advanced training in branches and staging. Syncthing \cite{syncthing} handles peer-to-peer transfers serverlessly, but it treats file version history as an afterthought and lacks a simple interface to undo accidental overwrites. 

Our goal was to build a tool that does all three: synchronize files automatically on a local network, protect files with automatic version backups before overwrites, and run with zero configuration or technical overhead.

\section{Context of the Project}

This project was developed as part of the Action Learning programme at EPITA, a programme that asks students to identify a real problem and build a solution in a team of four. Our team has two backend developers working in Go and C++, and two frontend and CI/CD developers working in Java with GitHub Actions.

The idea came from our own frustration. Rather than building something theoretical, we decided to solve the file-sharing problem we kept running into. The system we set out to build should feel invisible: put files into a folder on your laptop, and they show up on your teammates' laptops. Change a file, and everyone gets the update. If an update accidentally overwrites something, a backup of the old version is waiting in a local folder.

The Action Learning context shaped how we worked. With a fixed timeline and a small team, every design decision had to balance ambition against what could actually be built and tested in the available time.

\section{Objectives}

The project has three objectives:

\begin{enumerate}
    \item \textbf{Build a serverless synchronization system.} Peers discover each other automatically on a local network using mDNS/DNS-SD (the same protocols that let printers appear on your Wi-Fi without configuration). Files synchronize over direct TCP connections. No server, no accounts, no internet required.

    \item \textbf{Protect against accidental data loss.} Before any synchronized overwrite, the system saves the existing file into a local backup directory. Conflict resolution uses a Last-Write-Wins rule based on Lamport clock timestamps: predictable, automatic, and reversible.

    \item \textbf{Test two specific claims:}
    \begin{itemize}
        \item[$H_1$:] Average synchronization latency under one second on a standard WLAN for files up to 10\,MB, with no cloud infrastructure.
        \item[$H_2$:] For 3 to 5 peers, Lamport clocks with LWW are sufficient to maintain consistency and support rollback, without needing vector clocks or CRDTs.
    \end{itemize}
\end{enumerate}

\section{Overview of the Thesis}

Chapter~2 reviews the academic literature on distributed synchronization and collaborative file management. It discusses the trade-offs between different approaches and identifies the specific gap our project addresses.

Chapter~3 covers the state of the art: the research findings and technical decisions that determined our architecture, from protocol selection to conflict resolution strategy.

Chapter~4 describes the system architecture, how the Go and C++ components divide responsibilities, and how data flows through the system.

Chapter~5 explains the methodology, including our team structure, development process, tools, and the action plan for Action Learning Week.

Chapter~6 lists the expected development outcomes. After Action Learning Week, it will also cover what actually happened.

Chapter~7 presents the conclusions, summarises the contributions, and suggests directions for future work.

The appendices contain detailed comparison tables that support the analysis in Chapter~2.
